\documentclass[11pt]{amsart}
\usepackage[a4paper,margin=25mm]{geometry}
\usepackage{amssymb,microtype,booktabs}
\usepackage[hidelinks]{hyperref}
\allowdisplaybreaks
\emergencystretch=2em
\newtheorem{theorem}{Theorem}[section]\newtheorem{proposition}[theorem]{Proposition}\newtheorem{lemma}[theorem]{Lemma}
\newtheorem{corollary}[theorem]{Corollary}\newtheorem{conjecture}[theorem]{Conjecture}
\newtheorem{theoremA}{Theorem}\renewcommand{\thetheoremA}{\Alph{theoremA}}
\theoremstyle{definition}\newtheorem{definition}[theorem]{Definition}\newtheorem{example}[theorem]{Example}
\theoremstyle{remark}\newtheorem{remark}[theorem]{Remark}
\newcommand{\Q}{\mathbb Q}\newcommand{\Z}{\mathbb Z}\newcommand{\F}{\mathbb F}
\newcommand{\Gal}{\operatorname{Gal}}\newcommand{\disc}{\operatorname{disc}}\newcommand{\Res}{\operatorname{Res}}
\newcommand{\lead}{\operatorname{lc}}\newcommand{\ur}{\mathrm{ur}}\newcommand{\supp}{\operatorname{supp}}
\begin{document}
\title[Chebyshev--Lucas critical polynomials]{Chebyshev--Lucas critical polynomials:\\ local blocks, irreducibility and Galois groups}
\author{Alexandr Martinevski}
\address{Independent Researcher, Lithuania}\email{\href{mailto:alex.martinevski@gmail.com}{alex.martinevski@gmail.com}}
\thanks{ORCID: \href{https://orcid.org/0009-0000-4230-4414}{0009-0000-4230-4414}. Companion to \cite{I,A}.}
\date{September 2026}
\begin{abstract}
For $d\ge3$ let $A_d(x)=\sum_{k\ge0}(-1)^k\binom d{2k+1}x^k-1$, of degree $m=\lfloor(d-1)/2\rfloor$. Its roots are $\tan^2\theta$ at the solutions of
$\sin d\theta=\sin\theta\cos^{d-1}\theta$, the critical points of the extremal problem of \cite{I}, and $A_d+1$ is the Lucas polynomial $U_d(2,1+x)$.
We conjecture that $A_d$ is irreducible over $\Q$ for all $d$ and develop the tools for this problem. For even $d$ the $2$-adic Newton polygon
proves irreducibility whenever $\gcd(v_2(d),d/2-1)=1$ (about $94.5\%$ of even $d$). For odd $d$ we exclude rational roots, prove irreducibility
for $d=2p+1$ with $2$ a primitive root mod $p$, and show, via an explicit primitive divisor theorem, that factors of any fixed degree
disappear for large $d$. The central new tool is a local block lemma: for $\ell^a\Vert d-1$ the $\ell$-adic Newton polygon has vertices
$((\ell^b-1)/2,a-b)$, the reduction is $x^{(\ell^a-1)/2}H$, and $H$ is separable iff $\ell$ does not divide an explicit resultant $N_k$; for
even $d$ there is a dual block at every odd prime power $q^b\Vert d$, $b\ge1$, with no such condition. With these we compute Galois groups: $\Gal(A_{2p+1})\cong S_p$
when $2$ is a primitive root mod $p$; $\Gal(A_{2p+2})\supseteq A_p$ whenever $p+1\notin\{2^s,3\cdot2^s\}$; and
$\Gal(A_{2p+2})\cong S_p$ for every prime $p\equiv\pm3\pmod8$, an unconditional infinite family. Computations certify irreducibility for $d\le400$ and $\Gal(A_{2p+2})\cong S_p$ for all primes $p\le3000$. All roots of $A_d$ are real, and
irreducibility of $A_d$ determines the degree of the sharp constant: $[\Q(C_d):\Q]=d-1$ for odd $d$ and $d/2-1$ for even $d$. We explain why local methods alone cannot settle odd prime $d$.
\end{abstract}
\maketitle

\section{Introduction}\label{sec:intro}
\subsection*{The polynomials}
For an integer $d\ge3$ put $m=\lfloor(d-1)/2\rfloor$ and
\begin{equation}\label{eq:Ad}
A_d(x)=\sum_{k=0}^{m}(-1)^k\binom d{2k+1}x^k-1,\qquad Q_d(y)=y^mA_d\!\Bigl(\frac{1-y}{y}\Bigr).
\end{equation}
Thus $A_3=2-x$, $A_4=3-4x$, $A_5=x^2-10x+4$, $A_6=6x^2-20x+5$. With $x=\tan^2\theta$,
\begin{equation}\label{eq:trig}
A_d(\tan^2\theta)+1=\frac{\sin d\theta}{\sin\theta\,\cos^{d-1}\theta},
\end{equation}
so the roots of $A_d$ correspond to the solutions of $\sin d\theta=\sin\theta\cos^{d-1}\theta$. These are the critical points of
$\theta\mapsto\cos^d\theta-\cos d\theta$, whose maximum is the sharp constant $C_d=\|T_d-x^d\|_{\infty,[-1,1]}$ for multivariate
characteristic-function factorization \cite{I}. There $y=(c_d^*)^2$, with $c_d^*$ the maximizing cosine, is shown to be a root of an integer polynomial of degree $m$ with
leading coefficient $2^{d-1}-1$ \cite[Prop.~5.2]{I}; dividing the critical equation $U_{d-1}(c)=c^{d-1}$ by $c^{d-1}$ shows that this polynomial is exactly $Q_d$.
The polynomial $A_d$ appears in \cite[Remark~5.3]{I}, where its irreducibility for all $d$ is posed as an open problem \cite[Sec.~10]{I}. Irreducibility of $Q_d$ fixes the degree
of $y$; a separate field-degree argument then gives the degree of $C_d$. Equivalently, $A_d+1=U_d(2,1+x)$, where $U_n(P,Q)=(\alpha^n-\beta^n)/(\alpha-\beta)$, $\alpha+\beta=P$, $\alpha\beta=Q$,
is the Lucas polynomial. We call $A_d$ the \emph{Chebyshev--Lucas critical polynomials}.

\begin{conjecture}\label{conj:main}
$A_d$ is irreducible over $\Q$ for every $d\ge3$. Moreover $\Gal(A_d/\Q)\cong S_m$.
\end{conjecture}

The conjecture has the flavour of classical open irreducibility problems for binomial-type families (Cauchy--Mirimanoff polynomials,
truncated binomial expansions) and of recent work on Chebyshev-like sequences \cite{DKS}. We found no known theorem that covers $A_d$.

\subsection*{Main results}
Write $a=v_2(d)$ for even $d$. Our results fall into four groups.

\begin{theoremA}[even $d$]\label{thmA}
If $d$ is even, every irreducible factor of $A_d$ over $\Q_2$ has degree divisible by $m/\gcd(a,m)$. In particular $A_d$ is irreducible if
$\gcd(a,m)=1$, and also if $d-1$ is prime. Among even $d\le10^5$ these criteria cover $48015$ of $49999$ values.
\end{theoremA}

\begin{theoremA}[odd $d$]\label{thmB}
Let $d\ge5$ be odd. Then $A_d$ has no rational root. If $A_d$ has an irreducible factor of degree $k$, then
$d<\max\{2^{2k+1},10^{30}(2k)^9\}$. If $d=2p+1$ with $p$ prime and $2$ a primitive root modulo $p$, then $A_d$ is irreducible.
\end{theoremA}

\begin{theoremA}[local blocks]\label{thmC}
Let $\ell$ be an odd prime and $\ell^a\Vert d-1$. The $\ell$-adic Newton polygon of $A_d$ on the roots of positive valuation has vertices
$((\ell^b-1)/2,\,a-b)$, $b=0,\dots,a$; the first block is tamely and totally ramified of degree $(\ell-1)/2$ and the others are wild.
The reduction is $A_d\equiv x^{(\ell^a-1)/2}H\pmod\ell$ with explicit $H$, and $H$ is separable iff $\ell\nmid N_k$, $k=(d-1)/\ell^a$.
For even $d$ and an odd prime power $q^b\Vert d$, $b\ge1$, there is a dual block on the roots of negative valuation (the reversed polygon has
vertices $((q^c-1)/2,b-c)$), and the remaining reduction is always separable. Inertia at such primes acts as an explicit cycle, or product of cycles, and fixes all other roots.
\end{theoremA}

\begin{theoremA}[Galois groups]\label{thmD}
Let $p$ be a prime.
\begin{enumerate}
\item If $p\ge5$ and $2$ is a primitive root modulo $p$, then $\Gal(A_{2p+1}/\Q)\cong S_p$.
\item If $p+1\notin\{2^s,3\cdot2^s\}$, then $\Gal(A_{2p+2}/\Q)\supseteq A_p$.
\item If $p\equiv\pm3\pmod8$, then $\Gal(A_{2p+2}/\Q)\cong S_p$. By Dirichlet's theorem this family is infinite.
\end{enumerate}
\end{theoremA}

\begin{theoremA}[the degree of $C_d$]\label{thmE}
All roots of $A_d$ are real and positive. If $A_d$ is irreducible, then $[\Q(C_d):\Q]=d-1$ for odd $d$ and $[\Q(C_d):\Q]=d/2-1$ for even $d$.
This holds unconditionally for every $d$ covered by Theorems~\ref{thmA} and \ref{thmB} and for all $d\le400$.
\end{theoremA}
Since $\deg Q_d=m$, \cite[Prop.~5.2]{I} gives only $[\Q(y):\Q]\le m$; Theorem~\ref{thmE} determines the degree of $C_d$ itself for these $d$ and reduces it
to Conjecture~\ref{conj:main} in general.

In Theorem~\ref{thmD}, part (1) is unconditional for each $p$; the infinitude of such $p$ is Artin's conjecture for the base $2$, known under GRH \cite{Hooley}.
Part (3) needs no hypothesis at all. Theorem~\ref{thmD} is proved without the classification of finite simple groups (CFSG): besides the
Jordan--Marggraff theorem we use only Burnside's theorem on groups of prime degree and Bochert's bound on the minimal degree of doubly transitive
groups. The minimal degree theorem of Liebeck and Saxl \cite{LS}, which relies on CFSG, gives a stronger alternative (Remark~\ref{rem:LS}).

\subsection*{Computation}
Irreducibility is certified for all $3\le d\le400$ (Section~\ref{sec:comp}). Moreover $\Gal(A_d)\cong S_m$ for $17\le d\le82$,
$\Gal(A_{2p+1})\cong S_p$ for all primes $3\le p\le199$, and $\Gal(A_{2p+2})\cong S_p$ for all primes $3\le p\le3000$, the last with a
per-prime certificate table.

\subsection*{Limits of local methods}
For odd $d$ the question is the \emph{macroscopic} range of factor degrees. Periodicity of short cycles mod $p$ (Proposition~\ref{prop:period})
does not help: on explicit progressions of $d$ every fixed finite set of primes sees all short cycles (Proposition~\ref{prop:obstruction}).
For odd prime $d$ all local data at the structured primes are compatible with a factorization; see Example~\ref{ex:47} for $d=47$.
A proof for all $d$ needs a global argument.

\subsection*{Organization}
Section~\ref{sec:struct} collects identities. Section~\ref{sec:even} treats even $d$ and Section~\ref{sec:odd} odd $d$. Section~\ref{sec:blocks}
proves the local block lemmas and Section~\ref{sec:galois} the results on Galois groups. Section~\ref{sec:limits} discusses periodicity and the
limits of local methods, Section~\ref{sec:comp} the computations, and Section~\ref{sec:open} open problems.

\section{Structure}\label{sec:struct}
Put $a_k(x)=A_k(x)+1=\sum_i(-1)^i\binom k{2i+1}x^i$ and $b_k(x)=\sum_i(-1)^i\binom k{2i}x^i$, so that, with $x=-s^2$,
\begin{equation}\label{eq:ab}
(1+s)^k=b_k(x)+s\,a_k(x).
\end{equation}
The leading coefficient of $A_d$ is $(-1)^m$ for odd $d$ and $(-1)^md$ for even $d$; moreover $A_d(0)=d-1$ and $A_d(-1)=2^{d-1}-1$.

\begin{proposition}\label{prop:identities}
\begin{enumerate}
\item $2s\,(A_d(-s^2)+1)=(1+s)^d-(1-s)^d$ in $\Z[s]$.
\item $A_d(x)+1=U_d(2,1+x)=\sum_{j}(-1)^j\binom{d-1-j}{j}2^{d-1-2j}(1+x)^j$.
\item $\sum_{j=0}^{d-1}(2u-1)^j-u^{d-1}=u^{d-1}A_d\bigl(-(u-1)^2/u^2\bigr)$.
\item $A_d+1=c_d\prod_{e\mid d,\,e\ge3}\Psi_e$, where $\Psi_e\in\Z[x]$ is the minimal polynomial of $\tan^2(\pi/e)$, of degree $\varphi(e)/2$, and
$c_d\in\Q^\times$. For odd $d$ all $\tan^2(j\pi/d)$ are algebraic integers; for even $d$ the factor $\Psi_e$ is non-monic exactly when
$e=2q^c$ with $q$ an odd prime.
\end{enumerate}
\end{proposition}
\begin{proof}
(1) is \eqref{eq:ab} for $k=d$. For (2), in (1) put $\alpha=1+s$, $\beta=1-s$, so $\alpha+\beta=2$ and $\alpha\beta=1+x$. For (3), substitute
$s=(u-1)/u$ in (1): $1-s=1/u$ and $1+s=(2u-1)/u$. For (4), the roots of $A_d+1$ are $\tan^2(j\pi/d)$, $1\le j\le m$, by \eqref{eq:trig}; the
Galois orbits of these numbers are indexed by $e=d/\gcd(j,d)$. For the last statement write $\tan^2(\pi j/e)=-((\zeta-1)/(\zeta+1))^2$ with $\zeta$ a
primitive $e$-th root of unity. Then $1+\zeta$ is a unit unless $-\zeta$ has odd prime power order, i.e.\ unless $e=2q^c$.
\end{proof}

\begin{lemma}[real roots]\label{lem:real}
Put $\Phi(\theta)=A_d(\tan^2\theta)+1=\sin d\theta/(\sin\theta\cos^{d-1}\theta)$ on $(0,\pi/2)$. The equation $\Phi=1$ has exactly $m$ solutions, all simple:
one in $(\pi/(2d),\pi/d)$; one in each half of every positive lobe $(2i\pi/d,(2i+1)\pi/d)$ with $i\ge1$ and $(2i+1)\pi/d<\pi/2$; and, if $d\equiv1,2\pmod4$, one
in the first half of the last positive lobe, which reaches $\pi/2$. Hence $A_d$ has $m$ distinct positive roots $x_j=\tan^2\theta_j$, and $\Q(x_j)$
is totally real.
\end{lemma}
\begin{proof}
At the midpoint of a positive lobe $\sin d\theta=1>\sin\theta\cos^{d-1}\theta$, so $\Phi>1$ there, while $\Phi=0$ at the ends of the lobe. Hence each
half of a full positive lobe contains a solution. On $(0,\pi/d)$ we have $\Phi(0^+)=d$, $\Phi(\pi/(2d))>1$ and $\Phi(\pi/d)=0$. For $d=4k+1$ the last
lobe is $(2k\pi/d,\pi/2)$, its first half, with $\Phi\to+\infty$ at $\pi/2$. For $d=4k+2$ the last lobe ends at $\pi/2$, where again
$\Phi\to+\infty$. Counting over the four classes of $d$ modulo $4$ gives at least $1+2(k-1)+1=m$ ($d=4k+1,4k+2$), $1+2k=m$ ($d=4k+3$) and
$1+2(k-1)=m$ ($d=4k$) solutions in the stated intervals. Since $\deg A_d=m$, these are all the roots, and they are simple.
\end{proof}

\section{Even degree}\label{sec:even}
\begin{theorem}\label{thm:even}
Let $d$ be even, $a=v_2(d)$ and $m=d/2-1$. The $2$-adic Newton polygon of $A_d$ is the single segment from $(0,0)$ to $(m,a)$, and every
interior point lies strictly above it. Hence every irreducible factor of $A_d$ over $\Q_2$ has degree divisible by $m/\gcd(a,m)$, and $A_d$ is
irreducible over $\Q$ when $\gcd(a,m)=1$.
\end{theorem}
\begin{proof}
The constant term $d-1$ is odd and the leading coefficient is $\pm d$. For $1\le k<m$ the coefficient is $\pm\binom d{2k+1}$, and for odd $j$,
$\binom dj=\frac dj\binom{d-1}{j-1}$ has $2$-adic valuation at least $a>ak/m$. Dumas' theorem \cite{Dumas} gives the rest.
\end{proof}
If $a\ge2$ then $m$ is odd, so $\gcd(a,m)=1$ whenever $a$ is a power of $2$; this covers $d\equiv2\pmod4$ and $d\equiv4\pmod8$.

\begin{proposition}\label{prop:eis}
If $p$ is an odd prime, then $A_{p+1}$ is Eisenstein at $p$.
\end{proposition}
\begin{proof}
$A_{p+1}(0)=p$; for $3\le j\le p-2$, $\binom{p+1}j=\binom pj+\binom p{j-1}\equiv0\pmod p$; the leading coefficient is $\pm(p+1)$.
\end{proof}

\begin{definition}
The \emph{exceptional even set} is $\mathcal E=\{d \text{ even}: \gcd(v_2(d),d/2-1)>1 \text{ and } d-1\text{ not prime}\}$.
\end{definition}
It begins $56,248,296,320,344,352,392,512,\dots$ and has $1984$ elements up to $10^5$. The $2$-adic factorization is in fact explicit.

\begin{proposition}\label{prop:ore}
Let $d$ be even with $a=v_2(d)\ge2$, and put $g=\gcd(a,m)$ (odd, since $m$ is odd) and $Q=m/g$. Over $\Q_2$,
\[A_d=\prod_{\psi}F_\psi,\]
where $\psi$ runs over the irreducible factors of $y^g-1$ in $\F_2[y]$ and $F_\psi$ is irreducible of degree $Q\deg\psi$. Consequently every proper
factor of $A_d$ over $\Q$ has degree $Q\,s$, where $s$ is the sum of $\deg\psi$ over a proper nonempty set of these $\psi$. For $g=3$ the only possible
degrees are $Q$ and $2Q$; for $g=5$ they are $Q$ and $4Q$.
\end{proposition}
\begin{proof}
Apply Ore's theorem on the residual polynomial \cite{Ore,GMN} to $x^mA_d(1/x)$, whose Newton polygon is the single segment from $(0,a)$ to $(m,0)$
(Theorem~\ref{thm:even}). Its lattice points are at the abscissae $iQ$, $0\le i\le g$, and all interior coefficients lie strictly above the segment, so
the residual polynomial is $c_0+c_gy^g$ with $c_0\equiv d/2^a\equiv1$ and $c_g\equiv d-1\equiv1\pmod2$, i.e.\ $y^g+1$. It is separable because $g$
is odd, so $A_d$ is regular and its $\Q_2$-factors correspond to the irreducible factors of $y^g+1=y^g-1$ over $\F_2$, each of degree $Q$ times the
degree of the factor.
\end{proof}

\begin{remark}\label{rem:Estruct}
For $d\in\mathcal E$, the degrees allowed by Proposition~\ref{prop:ore} can be tested against the local data at the other structured primes: the
dual blocks and the separable reduction at each odd $q\mid d$ (Lemma~\ref{lem:dual}), and the blocks and $H$ at each odd $\ell\mid d-1$
(Lemma~\ref{lem:block}). For the $15$ elements of $\mathcal E$ below $1000$ this excludes every proper degree except for $d=248$ (degrees
$41,82$) and $d=352$ (degrees $35,140$), which are settled by the general certificates of Section~\ref{sec:comp}. The exclusions depend on the
factorization patterns of these reductions, so they give no uniform argument for $\mathcal E$.
\end{remark} For $d=\ell^c+1$ with $c\ge2$, Lemma~\ref{lem:block}
shows that the $\ell$-adic polygon is never a single segment, so Proposition~\ref{prop:eis} does not extend to prime powers.

\section{Odd degree}\label{sec:odd}
\begin{lemma}\label{lem:mod2}
For odd $d$, $A_d(x-1)\equiv x^m-1\pmod 2$.
\end{lemma}
\begin{proof}
Modulo $2$ we have $(2u-1)^j\equiv1$, so Proposition~\ref{prop:identities}(3) gives $u^{2m}\bar A(-(u-1)^2/u^2)=1+u^{2m}$ in $\F_2[u]$, where
$\bar A=A_d\bmod 2$. Since $(u+1)^{2k}=(u^2+1)^k$ in characteristic $2$, with $U=u^2$ this reads $U^m\bar A(1+1/U)=1+U^m$, i.e.
$\bar A(1+Y)=1+Y^m$.
\end{proof}

\begin{lemma}\label{lem:units}
Let $d$ be odd, $x$ a root of $A_d$, $s^2=-x$, $\alpha=1+s$, $\beta=1-s$, $\gamma=\alpha/\beta$, $K=\Q(s)$. Then $\alpha,\beta$ are units at every
prime of $K$ above $2$, and $\gamma$ is not a root of unity.
\end{lemma}
\begin{proof}
All these numbers are algebraic integers, because $A_d$ is monic up to sign. Let $\mathfrak P\mid2$. By Lemma~\ref{lem:mod2},
$A_d(t-1)=t^m-1+2R(t)$ with $R\in\Z[t]$, so $t=1+x=\alpha\beta$ satisfies $t^m\equiv1\pmod{\mathfrak P}$ and $\alpha,\beta$ are $\mathfrak P$-units.
Hence $\gamma+1=2/\beta$ has $\nu_{\mathfrak P}(\gamma+1)=\nu_{\mathfrak P}(2)>0$. Suppose $\gamma$ is a root of unity. Then $\gamma\ne-1$, since
$\alpha+\beta=2$, so $\eta=-\gamma\ne1$ is a root of unity with $\nu(1-\eta)=\nu(2)$. If the order of $\eta$ is not a power of $2$, then $1-\eta$
is a $2$-unit. If the order is $2^c$, then $\nu(1-\eta)=\nu(2)/2^{c-1}$, which equals $\nu(2)$ only for $c=1$, i.e.\ $\gamma=1$ and $x=0$. But
$A_d(0)=d-1\ne0$.
\end{proof}

\begin{proposition}\label{prop:norational}
For odd $d\ge5$, $A_d$ has no rational root.
\end{proposition}
\begin{proof}
A rational root is an integer $r\mid d-1$. For $z\ge0$, $A_d(-z)=\sum_k\binom d{2k+1}z^k-1\ge d-1>0$, so $r>0$. By Lemma~\ref{lem:mod2},
$A_d(r)\equiv(r+1)^m+1\pmod2$, so $r$ is even. Then $(1\pm\sqrt{-r})$ is a Lucas pair: $\alpha+\beta=2$ and $\alpha\beta=1+r$ are coprime, and
$\alpha/\beta$ is not a root of unity by Lemma~\ref{lem:units}. By Proposition~\ref{prop:identities}(1), $A_d(r)=0$ iff $u_d(\alpha,\beta)=1$. For
$d>30$ the number $u_d$ has a primitive divisor \cite[Thm~1.4]{BHV}, so $|u_d|>1$. For $5\le d\le30$ one checks the divisors of $d-1$ directly.
\end{proof}

\begin{theorem}\label{thm:2p1}
Let $p$ be an odd prime with $\operatorname{ord}_p(2)=p-1$ and $d=2p+1$. Then $A_d$ is irreducible.
\end{theorem}
\begin{proof}
Here $m=p$ and $A_d(x-1)\equiv(x-1)\Phi_p(x)\pmod2$, where $\Phi_p$ is irreducible over $\F_2$ because $2$ has order $p-1$. The leading
coefficient is odd, so a proper factorization over $\Z$ would have a factor of degree $1$, i.e.\ a rational root. This contradicts
Proposition~\ref{prop:norational}.
\end{proof}

\begin{theorem}[small factors]\label{thm:small}
Let $d$ be odd. If $A_d$ has an irreducible factor of degree $k$ over $\Q$, then $d<\max\{2^{2k+1},10^{30}(2k)^9\}$.
\end{theorem}
\begin{proof}
Let $x$ be a root of such a factor, with $s,\alpha,\beta,\gamma$ as in Lemma~\ref{lem:units}. Since $s=(\gamma-1)/(\gamma+1)$, $\Q(\gamma)=\Q(s)$
has degree $D\le2k$, and $\beta=2/(\gamma+1)$ is an algebraic integer of $\Q(\gamma)$. Also $x\ne0,-1$, so $\gamma\ne0,1$. By
Proposition~\ref{prop:identities}(1), $\alpha^d-\beta^d=\alpha-\beta$, i.e.\ $\gamma^d-1=(\gamma-1)\beta^{1-d}$. Hence for every prime
$\mathfrak p$ of $\Q(\gamma)$,
\[\nu_{\mathfrak p}(\gamma^d-1)=\nu_{\mathfrak p}(\gamma-1)+(1-d)\nu_{\mathfrak p}(\beta)\le\nu_{\mathfrak p}(\gamma-1),\]
so $\gamma^d-1$ has no primitive divisor. Since $\gamma$ is not a root of unity (Lemma~\ref{lem:units}), the explicit form \cite[Thm~4.2]{BL} of Schinzel's
theorem \cite{Schinzel} gives $d<\max\{2^{D+1},10^{30}D^9\}$.
\end{proof}
\begin{remark}
The argument uses that the constant is $-1$. For $U_d(2,t)-c$ one gets $\gamma^d-1=c(\gamma-1)\beta^{1-d}$, and a primitive divisor may divide
$c$. Numerically, $U_d(2,t)-c$ is irreducible for $c\in\{\pm1,\pm2,3,4\}$ and $d\le50$.
\end{remark}

\section{Local blocks}\label{sec:blocks}
Put $N_k=\prod_{\zeta^k=1,\,\zeta\ne1}\bigl((2-\zeta)^k-1\bigr)\in\Z\setminus\{0\}$; it is nonzero because $|2-\zeta|>1$ for $\zeta\ne1$ on the
unit circle. Equivalently, $N_k=\Res\bigl((x^k-1)/(x-1),\,(2-x)^k-1\bigr)$, and $\Res(R_k,S_k)=\pm N_k/k$ for $R_k(z)=((1+z)^k-1)/z$ and
$S_k(z)=((1-z)^k-1)/z$.

\begin{lemma}[block at $\ell^a\Vert d-1$]\label{lem:block}
Let $\ell$ be an odd prime, $d=k\ell^a+1$ with $\ell\nmid k$, $N=\ell^a$, $E=(N-1)/2$ and $E_b=(\ell^b-1)/2$.
\begin{enumerate}
\item[(A)] $v_\ell\binom d{2i+1}=a-v_\ell\bigl(i(2i+1)\bigr)$ for $1\le i\le E$.
\item[(B)] The $\ell$-adic Newton polygon of $A_d$ on its positive-valuation part has vertices $(E_b,a-b)$, $b=0,\dots,a$. Its segments have height
$1$ and lengths $\ell^{b-1}(\ell-1)/2$. The first block is tamely and totally ramified of degree $e=(\ell-1)/2$; for $b\ge2$ the blocks are wild.
\item[(C)] $A_d\equiv x^EH\pmod\ell$, with $H=(-1)^Ea_k^N+x^{E+1}w^N$, $w=(b_k-1)/x$ and $H(0)\equiv(-1)^Ek\not\equiv0$.
\item[(D)] $H$ is separable iff $\ell\nmid N_k$, iff there are no $\zeta_1\ne\zeta_2$ in $\mu_k(\overline{\F}_\ell)$ with $\zeta_1+\zeta_2=2$.
\end{enumerate}
\end{lemma}
\begin{proof}
(A) We have $\binom dj=\frac dj\binom{d-1}{j-1}$ and $v_\ell(d)=0$. For $1\le t<N$, $v_\ell\binom{kN}t=a-v_\ell(t)$, since
$\binom{kN-1}{t-1}\not\equiv0$ by Lucas' theorem. Finally $j=2i+1$ and $j-1=2i$ are not both divisible by $\ell$.
(B) A point with $v_\ell(i(2i+1))=t$ has $i\ge E_t$, so it lies on or to the right of the vertex $(E_t,a-t)$ at the same height. The slopes
$-2/(\ell^{b-1}(\ell-1))$ increase. Each segment has height $1$, so each block is irreducible over $\Q_\ell^\ur$ and totally ramified, of degree
$\ell^{b-1}e$.
(C) In $\F_\ell[s]$, $(1\pm s)^d=(1\pm s)(1\pm s^N)^k$. By \eqref{eq:ab} and $a_k(x^N)=a_k(x)^N$,
$A_d+1\equiv(-x)^Ea_k^N+b_k^N$, and $b_k-1=xw$.
(D) Modulo $\ell$, $H'\equiv(E+1)x^Ew^N$ with $E+1=(N+1)/2\not\equiv0$. A repeated root $r$ of $H$ has $r\ne0$, hence $w(r)=0$ and then
$a_k(r)=0$. With $s^2=-r$ this means $(1\pm s)^k=1$; conversely such $s$ gives a repeated root. Reducing $\mu_k$ modulo a prime above $\ell$
gives the resultant form.
\end{proof}

\begin{lemma}[tame subfield]\label{lem:contain}
In Lemma~\ref{lem:block}, let $K_0=\Q_\ell^\ur$ and $T=K_0(\ell^{1/e})$. Then $K_0(\beta)\supseteq T$ for every root $\beta$ of every block.
\end{lemma}
\begin{proof}
$K_0(\beta)/K_0$ is totally ramified of degree $\ell^{b-1}e$. Its maximal tamely ramified subextension has degree equal to the prime-to-$\ell$
part of the ramification index, namely $e$ \cite[Ch.~IV]{Serre}. Since the residue field of $K_0$ is algebraically closed and $\ell\nmid e$,
$T$ is the only tamely ramified extension of $K_0$ of degree $e$.
\end{proof}

\begin{corollary}[inertia at $\ell$]\label{cor:inertia}
In Lemma~\ref{lem:block}:
\begin{enumerate}
\item if $a=1$ and $\ell\nmid N_k$, inertia at $\ell$ contains an $e$-cycle that fixes every other root;
\item if $a\ge2$ and $\ell\nmid N_k$, an element of order $e$ that generates the tame part of inertia is a product of $(\ell^a-1)/(\ell-1)$ disjoint
$e$-cycles on the block roots and fixes the other roots; no inertia element is a single $e$-cycle;
\item inertia at $\ell$ contains an odd permutation iff $\ell\equiv1\pmod4$ and $a$ is odd.
\end{enumerate}
\end{corollary}
\begin{proof}
If $\ell\nmid N_k$, then by Hensel the roots outside the blocks lie in $K_0$ and are fixed by inertia. Now $T/K_0$ is cyclic of degree $e$ and
acts simply transitively on the roots of the first block. By Lemma~\ref{lem:contain}, an element that is nontrivial on $T$ fixes no root of any
block, which gives (1) and (2). For (3), wild inertia has odd order $\ell^r$, and the number of $e$-cycles is $\equiv a\pmod2$.
\end{proof}

\begin{lemma}[dual block]\label{lem:dual}
Let $d$ be even, $q$ an odd prime, $q^b\Vert d$, $N=q^b$, $E=(N-1)/2$ and $e=d/N$.
\begin{enumerate}
\item $A_d\equiv(-x)^Ea_e(x)^N-1\pmod q$, a polynomial of degree $m-E$.
\item This reduction is separable.
\item The $q$-adic polygon of $x^mA_d(1/x)$ has vertices $((q^c-1)/2,\,b-c)$, $c=0,\dots,b$, on its positive-slope part. Thus $A_d$ has exactly
$E$ roots of negative valuation, forming blocks of degrees $q^{c-1}(q-1)/2$.
\end{enumerate}
Consequently inertia at $q$ fixes all roots of non-negative valuation; for $b=1$ and $q\ge5$ it contains an $E$-cycle fixing all other roots;
for $b\ge2$ its elements move at most $E$ roots; and it contains an odd permutation iff $q\equiv1\pmod4$ and $b$ is odd.
\end{lemma}
\begin{proof}
(1) In $\F_q[s]$, $(1\pm s)^d=(1\pm s^N)^e=b_e(x)^N\pm s^Na_e(x)^N$. Take the difference and divide by $2s$. (2) The derivative is
$\equiv-E(-x)^{E-1}a_e^N$, which is nonzero at every root, since a root has $x\ne0$ and $a_e(x)\ne0$. (3) The coefficient of $x^j$ in the
reversed polynomial is $\pm\binom d{2j+1}$, and $v_q\binom{eN}t=b-v_q(t)$ for $1\le t<N$ while $v_q\binom{eN}N=0$; argue as in
Lemma~\ref{lem:block}(B). The inertia statements follow from Lemma~\ref{lem:contain} and Corollary~\ref{cor:inertia} applied to $y=1/x$.
\end{proof}
For odd $d$ the leading coefficient is $\pm1$; the reduction (1) then has full degree and is separable, so $q\mid d$ is unramified.

\begin{lemma}[support bound]\label{lem:support}
Let $\ell$ be prime. Take $f=A_d$ if $\ell$ does not divide its leading coefficient and $f=x^mA_d(1/x)$ otherwise. Write
$\bar f=c\prod_i\phi_i^{e_i}$ over $\F_\ell$ and put $M_\ell(d)=\sum_{e_i\ge2}e_i\deg\phi_i$. Then every element of every inertia group above
$\ell$ moves at most $M_\ell(d)$ roots of $A_d$.
\end{lemma}
\begin{proof}
A simple root of $\bar f$ lifts by Hensel to exactly one root of $f$ in an unramified extension; these $m-M_\ell$ roots are fixed by inertia.
\end{proof}

\section{Galois groups}\label{sec:galois}
We use the following facts about a transitive group $G\le S_n$.
\begin{enumerate}
\item[(a)] If $n$ is prime, $G$ is primitive.
\item[(b)] (Jordan--Marggraff) If $G$ is primitive and contains a subgroup with an orbit of length $k$, $1<k<n/2$, fixing the remaining points,
then $G\supseteq A_n$ \cite[Thms~13.4--13.5]{Wielandt}; see \cite{LT} for a short proof of Marggraff's theorem and \cite{Jones} for its history.
\item[(c)] (Jones) If $G$ is primitive and contains a cycle fixing at least three points, then $G\supseteq A_n$ \cite[Cor.~1.3]{Jones}.
\item[(d)] (Burnside--Bochert) If $n=p\ge41$ is prime and $G$ contains $g\ne1$ moving fewer than $p/4$ points, then $G\supseteq A_p$.
\end{enumerate}
For (d): by Burnside's theorem a transitive group of prime degree $p$ that is not doubly transitive is a subgroup of $\mathrm{AGL}(1,p)$ \cite{Muller}.
A nonidentity element of $\mathrm{AGL}(1,p)$ moves at least $p-1$ points, so $g$ excludes this case. By Schomburg's classification-free form of
Bochert's theorem \cite[Thm~3.1]{Schomburg} (see also \cite{Bochert97}), a doubly transitive group of degree $n\ge38$ not containing $A_n$ has
minimal degree at least $n/4$. Items (a), (b), (d) do not use CFSG; (c) does, and it is not needed in the proofs below.

\begin{remark}[Liebeck--Saxl]\label{rem:LS}
Using CFSG, a primitive group of minimal degree $<n/3$ has $n=\binom{m'}{k}^r$ and $A_{m'}^r\le G\le S_{m'}\wr S_r$ in the product action on
$k$-sets \cite{LS}. For $n=p$ prime this forces $r=1$ and $k\in\{1,m'-1\}$. Hence in (d) one may replace $p/4$ by $p/3$ and drop $p\ge41$. This is
needed only in one subcase of Theorem~\ref{thm:Geven}(2), which Theorem~\ref{thmD} does not use.
\end{remark}

\begin{theorem}\label{thm:Gstar}
Let $p\ge5$ be prime, $d=2p+1$, and suppose $A_d$ is irreducible. Then $\Gal(A_d)\supseteq A_p$. If moreover $p\equiv1\pmod4$ or
$p\equiv3\pmod8$, then $\Gal(A_d)\cong S_p$.
\end{theorem}
\begin{proof}
$G$ is transitive of prime degree $p$, hence primitive. Apply Lemma~\ref{lem:block} with $\ell=p$, $a=1$, $k=2$: here $N_2=8$, so the block
is clean for every odd $p$. (Directly, $A_d\equiv(-x)^e((-x)^{e+1}+2)\pmod p$ with $e=(p-1)/2$.) Corollary~\ref{cor:inertia}(1) gives an
$e$-cycle, $1<e<p/2$, fixing $e+1\ge3$ roots, and (b) gives $A_p\subseteq G$. If $p\equiv1\pmod4$, the $e$-cycle is odd. If $p\equiv3\pmod8$,
then $(2/p)=-1$, so $f=\operatorname{ord}_p(2)$ is even and $(p-1)/f$ is odd. The Frobenius at $2$ (Lemma~\ref{lem:mod2}, $m=p$ odd) has one fixed
point and $(p-1)/f$ cycles of length $f$, so it is odd.
\end{proof}
\begin{corollary}\label{cor:G}
If $p\ge5$ is prime and $\operatorname{ord}_p(2)=p-1$, then $\Gal(A_{2p+1}/\Q)\cong S_p$.
\end{corollary}
\begin{proof}
Irreducibility is Theorem~\ref{thm:2p1}. The Frobenius at $2$ has cycle type $(1,p-1)$, which is odd.
\end{proof}

\begin{theorem}\label{thm:Geven}
Let $p\ge5$ be prime and $d=2p+2$ (so $m=p$); $A_d$ is irreducible by Theorem~\ref{thm:even}.
\begin{enumerate}
\item If $p+1\notin\{2^s,3\cdot2^s\}$, then $\Gal(A_d)\supseteq A_p$.
\item If some $\ell^a\Vert2p+1$ has $\ell^a\ge5$, $k=(2p+1)/\ell^a\ge3$ and $\ell\nmid N_k$, then $\Gal(A_d)\supseteq A_p$.
\end{enumerate}
\end{theorem}
\begin{proof}
(1) Suppose $p+1$ has a prime $q\ge5$ to the first power. Then $q\Vert d$, and Lemma~\ref{lem:dual} gives an $E$-cycle, $2\le E=(q-1)/2<p/2$,
fixing at least three roots, so (b) applies. Otherwise some odd $q^b\Vert p+1$ has $b\ge2$. Its inertia is nontrivial, since the polygon has a segment of non-integral slope, and moves at most
$E=(q^b-1)/2\le(p-1)/4<p/4$ roots. For $p\ge41$, (d) applies. The only prime $p<41$ in this case is $p=17$; there $2p+1=35$, and part (2) applies with
$\ell=7$, $a=1$, $k=5$, $N_5\not\equiv0\pmod7$ (Jordan--Marggraff).
(2) If $a=1$, Corollary~\ref{cor:inertia}(1) gives an $e$-cycle with $2\le e<p/2$ fixing at least three roots. If $a\ge2$, inertia is nontrivial and moves at most $E=(\ell^a-1)/2$ roots by Lemma~\ref{lem:support}. Now $E<p/4$ iff $k\ge5$, so (d) applies
when $k\ge5$ and $p\ge41$. For $k=3$ we have $E=(p-1)/3<p/3$, and Remark~\ref{rem:LS} applies (this subcase uses CFSG).
\end{proof}

\begin{theorem}[$2$-adic square class]\label{thm:T2}
Let $d$ be even with $a=v_2(d)\ge2$ and $\gcd(a,m)=1$. Then $\disc(A_d)\in(\Q_2^\times)^2$ iff $m\equiv\pm1\pmod8$.
\end{theorem}
\begin{proof}
By Theorem~\ref{thm:even} all roots have valuation $-a/m$, so $K=\Q_2(\text{root})$ is totally and tamely ramified of odd degree $m$. Such a field
is $\Q_2(\varpi^{1/m})$ for a uniformizer $\varpi=2u$ \cite[Ch.~IV]{Serre}. Every unit of $\Z_2$ is an $m$-th power: this is clear for $\pm1$, and on
$1+4\Z_2\cong\Z_2$ the map $v\mapsto mv$ is bijective. Hence $K=\Q_2(2^{1/m})$. The reversed polynomial has unit leading coefficient $d-1$ and the
same discriminant, so modulo squares $\disc(A_d)\equiv\disc(\mathcal O_K)\equiv\disc(x^m-2)=(-1)^{(m-1)/2}m^m2^{m-1}\equiv(-1)^{(m-1)/2}m$. This is
a $2$-adic square iff it is $\equiv1\pmod8$, i.e.\ iff $m\equiv\pm1\pmod8$.
\end{proof}

\begin{corollary}\label{cor:Sp}
For every prime $p\equiv\pm3\pmod8$, $\Gal(A_{2p+2}/\Q)\cong S_p$.
\end{corollary}
\begin{proof}
Here $v_2(2p+2)\ge2$. If $p\equiv5\pmod8$ and $p\ge13$, then $p+1=2\cdot o$ with $o\ge7$ odd. If $p\equiv3\pmod8$ and $p\ge19$, then
$p+1=4\cdot o$ with $o\ge5$ odd. In both cases $p+1\notin\{2^s,3\cdot2^s\}$, and Theorem~\ref{thm:Geven}(1) gives $A_p\subseteq G$.
Theorem~\ref{thm:T2} with $m=p\equiv\pm3$ gives a non-square discriminant, so $G=S_p$. The remaining primes $p=3,5,11$ are covered by
Section~\ref{sec:comp}.
\end{proof}

For the other residues the parity comes from closed-form local square classes. The decomposition group at $q$ lies in $A_m$ iff
$\disc\in(\Q_q^\times)^2$. Inertia alone does not detect this: for $d=64$ and $d=76$ the odd valuations $v_3=7$ and $v_5=17$ come from wild
blocks. Let $\chi_q$ be the Legendre symbol.

\begin{proposition}[dual square class]\label{prop:PD}
In Lemma~\ref{lem:dual} let $b=1$, and put $n=m-E$, $\delta=e/2-1$, $L=(-1)^\delta e$ and
$\bar D=-(-1)^{n(n-1)/2+(n+1)\delta}(-E)^nL^{n-1}$. Then $\disc(A_d)\notin(\Q_q^\times)^2$ iff $E$ is even, or $E$ is odd and
$\chi_q\bigl((-1)^{(E-1)/2}E\bigr)\chi_q(\bar D)=-1$.
\end{proposition}
\begin{proof}
The reversed polynomial $R$ has unit leading coefficient and $\disc R=\disc A_d$. Over $\Z_q$, $R=(d-1)R_BR_C$, where $R_B$ is Eisenstein of degree
$E$ and $R_C$ is monic with $\bar R_C$ proportional to the reversal of $\bar A=(-x)^Ea_e^q-1$. Modulo squares, $\disc R\equiv\disc R_B\disc R_C$.

For the block, $\disc R_B=(-1)^{E(E-1)/2}N(R_B'(\beta))$ and $R_B'(\beta)=E\beta^{E-1}(1+\varepsilon)$ with $v(\varepsilon)>0$. The norm of $1+\varepsilon$
is a square, so the class is $(-1)^{E(E-1)/2}E^EN(\beta)^{E-1}$. Since $N(\beta)$ has valuation $1$, this has odd valuation for even $E$ and
class $(-1)^{(E-1)/2}E$ for odd $E$.

For the rest, $\disc R_C$ is a unit congruent, modulo squares, to $\disc\bar A$. In $\F_q[x]$ we have $\bar A'=-E(-x)^{E-1}a_e^q$, of degree $n-1$,
and $\bar A=-1$ at the roots of $a_e$. So $\Res(\bar A,\bar A')=(-E)^n(-1)^{E-1}(-1)^{(n+1)\delta}L^n$ and $\lead\bar A=(-1)^EL$, which gives
$\disc\bar A=\bar D$.
\end{proof}

\begin{proposition}[square class at $\ell\Vert d-1$]\label{prop:PL}
Let $d$ be even, $\ell\Vert d-1$, $k=(d-1)/\ell\ge3$ and $\ell\nmid N_k$. Put $e=(\ell-1)/2$, $n=\ell(k-1)/2$, $\delta_a=(k-1)/2$,
$\delta_w=(k-3)/2$, $\lambda=(-1)^{\delta_a}k$, $\lambda_H=(-1)^{e+\delta_a}$ and $\rho=\Res(w,a_k)\bmod\ell\ne0$. Let
\[\bar D_H=(-1)^{n(n-1)/2}\lambda_H^{\,e-1}(e+1)^n\bigl((-1)^{n+e}k\bigr)^e(-1)^{(n+e)\delta_w}\lambda^{n-\delta_a}\rho .\]
Then $\disc(A_d)\notin(\Q_\ell^\times)^2$ iff $e$ is even, or $e$ is odd and $\chi_\ell\bigl((-1)^{(e-1)/2}e\bigr)\chi_\ell(\bar D_H)=-1$.
\end{proposition}
\begin{proof}
The argument is the same as for Proposition~\ref{prop:PD}, with the Eisenstein block of degree $e$ at $0$ and $H$ from Lemma~\ref{lem:block}(C).
Here $H'\equiv(e+1)x^ew^\ell$ has degree $n-1-e$, so $\disc H=(-1)^{n(n-1)/2}\lead(H)^{e-1}\Res(H,H')$. Moreover
\[\Res(H,x)=(-1)^nH(0),\qquad \Res(H,w)=(-1)^{n\delta_w}\lambda^n\!\!\prod_{w(r)=0}\!\!H(r),\]
with $H(r)=(-1)^ea_k(r)^\ell$ at the roots of $w$ and $\prod_{w(r)=0}a_k(r)=\rho\lambda^{-\delta_a}$.
\end{proof}
Theorem~\ref{thm:T2} and Propositions~\ref{prop:PD}, \ref{prop:PL} agree with the square class of the exact discriminant in $63$, $144$ and $125$ of
$63$, $144$ and $125$ cases respectively (all admissible even $d\le290$ in the tested ranges).

\section{Periodicity and the limits of local methods}\label{sec:limits}
\begin{proposition}[periodicity]\label{prop:period}
Let $p$ be an odd prime and $\ell\ge1$. The set of nonzero roots of $A_d\bmod p$ in $\F_{p^\ell}$ depends only on $d\bmod(p^{2\ell}-1)$, and
$x=0$ is a root iff $p\mid d-1$. Hence the number of distinct monic irreducible factors of degree $\ell$, other than $x$, is periodic in $d$ with
period dividing $p^{2\ell}-1$.
\end{proposition}
\begin{proof}
For $x\in\F_{p^\ell}^\times$ choose $s\in\F_{p^{2\ell}}$ with $s^2=-x$. By Proposition~\ref{prop:identities}(1), $A_d(x)=0$ iff
$(1+s)^d-(1-s)^d=2s$, which depends only on $d\bmod(p^{2\ell}-1)$. For $s=\pm1$ the condition is $2^{d-1}=1$, which depends on $d$ modulo a
divisor of $p-1$.
\end{proof}

\begin{proposition}[obstruction]\label{prop:obstruction}
If $d\equiv1\pmod{p^{2\ell}-1}$, then every $x\in\F_{p^\ell}^\times$ is a root of $A_d\bmod p$. Hence, for any finite set $P$ of odd primes, on the
progression $d\equiv1\pmod{\operatorname{lcm}_{p\in P}p(p^{2\ell}-1)}$ the factorization pattern at every $p\in P$ contains all cycles of length
dividing $\ell$; for odd $d$ the reduction mod $2$ also has a root. Cycle-type data at finitely many primes can therefore never exclude linear
factors for all $d$, and small factors need global input such as Proposition~\ref{prop:norational} and Theorem~\ref{thm:small}.
\end{proposition}
\begin{proof}
For $s\in\F_{p^{2\ell}}$ we have $(1\pm s)^d=1\pm s$.
\end{proof}

\begin{example}[$d=47$]\label{ex:47}
Here $m=23$, and the structured primes are $2$ and the prime divisors of $d(d-1)=47\cdot46$. Modulo $2$ the Frobenius has cycle type $(1,11,11)$,
since $2$ has order $11$ modulo $23$. At $23$ there is a tame Eisenstein block of degree $11$ and an unramified rest of degree $12$. At $47=d$,
$A_d\equiv(-x)^{23}-1$ splits completely, so the Frobenius is trivial. All these data are compatible with a factorization of degrees $11+12$ in
which the degree-$11$ factor is the $23$-adic block. Hence no argument that uses only these primes can prove irreducibility here (it holds by
computation). The same holds for $d=83,383,479$. For the composite $d=2p+1$ with $\operatorname{ord}_p(2)=(p-1)/2$ and $d\le400$, the Frobenius at
the primes $q\mid d$ closes every case except $d=15$.
\end{example}

\begin{remark}[orbit norms]
Let $d$ be odd and $A_d=FG$ in $\Z[x]$. Proposition~\ref{prop:identities}(4) gives $\prod_{x^0\in O_e}F(x^0)=\pm1$ for every orbit $O_e$ of
$\tan^2(j\pi/d)$, because $A_d(x^0)=-1$. This is the algebraic-integer analogue of Schur's argument for $\prod(x-a_i)-1$. It is, however,
automatic (it equals $\Res(\Psi_e,F)\Res(\Psi_e,G)=\pm1$), so it is only a necessary condition on the root set of a putative factor. For prime
$d$ it is a single real equation.
\end{remark}

\section{The degree of the sharp constant}\label{sec:degree}
Let $g(\theta)=\cos^d\theta-\cos d\theta$. By \cite[Thm~2.4]{I}, $C_d=\max_{[0,\pi/d]}g$ is attained at a unique interior critical point; since
$g'=d(\sin d\theta-\sin\theta\cos^{d-1}\theta)$, the critical points in $(0,\pi/2)$ are the $\theta_j$ of Lemma~\ref{lem:real}, and $C_d=g(\theta_1)$.
Put $c_j=\cos\theta_j$ and $y_j=c_j^2$; the $y_j$ are the roots of $Q_d$.

\begin{lemma}\label{lem:gvalues}
\begin{enumerate}
\item $g(\theta_j)=U_{d-2}(c_j)$.
\item If $\cos d\theta_j<0$ then $1<g(\theta_j)<2$; if $\cos d\theta_j>0$ then $-1<g(\theta_j)<0$.
\item Along the roots with $\cos d\theta_j<0$ the values $g(\theta_j)$ strictly decrease; along those with $\cos d\theta_j>0$ the values
$|g(\theta_j)|$ strictly increase.
\end{enumerate}
In particular the numbers $|g(\theta_j)|$, $1\le j\le m$, are nonzero and pairwise distinct.
\end{lemma}
\begin{proof}
(1) $\sin d\theta=\sin\theta\,U_{d-1}(\cos\theta)$, so at a critical point $U_{d-1}(c)=c^{d-1}$; with $T_d=cU_{d-1}-U_{d-2}$ we get
$g=c^d-T_d(c)=U_{d-2}(c)$.
(2) Put $u=\cos^d\theta$, $t=\tan\theta$ and $s=\sin\theta\cos^{d-1}\theta=ut=\sin d\theta$, so $\cos d\theta=\pm\sqrt{1-s^2}$. With $x=t^2$ we have
$ut^2=x(1+x)^{-d/2}\le M_d:=\frac2{d-2}(1-\frac2d)^{d/2}\le\frac2{e(d-2)}<1$. If $\cos d\theta<0$ then $g-1=u-s^2/(1+\sqrt{1-s^2})\ge u(1-M_d)>0$.
If $\cos d\theta>0$ then $g<0$ because $u^2+s^2=\cos^{2d-2}\theta<1$.
(3) By Lemma~\ref{lem:real}, consecutive roots of the same type satisfy $\theta'-\theta>h:=3\pi/(2d)$. Since
$\cos\theta'\le\cos\theta\cos(\theta'-\theta)$, the value $u'=\cos^d\theta'$ satisfies $u'/u\le\cos^dh\le e^{-dh^2/2}=e^{-9\pi^2/(8d)}$. For the first
type, $g-1\in[u(1-M_d),u]$, and for the second, $1-|g|\in[u,u(1+M_d)]$. So both claims follow from
\[e^{-9\pi^2/(8d)}<1-\frac2{e(d-2)}\qquad(d\ge3),\]
which also gives $e^{-9\pi^2/(8d)}(1+M_d)<1$. For $d\ge20$ it follows from $e^{-a}\le1-a+a^2/2$ with $a=9\pi^2/(8d)$; for $3\le d\le19$ it is checked
directly (the margin at $d=3$ is $0.239$).
\end{proof}

\begin{theorem}\label{thm:degree}
If $A_d$ is irreducible, then $[\Q(C_d):\Q]=d-1$ for odd $d$ and $[\Q(C_d):\Q]=d/2-1$ for even $d$.
\end{theorem}
\begin{proof}
$Q_d$ is irreducible of degree $m$, since $A_d$ is. For even $d$, $U_{d-2}$ is even, so $C_d=S(y_1)$ with $S\in\Q[y]$. The embeddings of $\Q(y_1)$ send
$C_d$ to $S(y_j)=g(\theta_j)$, which are distinct by Lemma~\ref{lem:gvalues}; hence the degree is $m=d/2-1$.

For odd $d$, $U_{d-2}$ is odd. The norm of $y_1$ is $\prod_jy_j=|Q_d(0)|/(2^{d-1}-1)=1/(4^m-1)$. This is not the square of a rational number,
because $4^m-1$ lies strictly between $(2^m-1)^2$ and $(2^m)^2$. So $c_1\notin\Q(y_1)$, $[\Q(c_1):\Q]=2m$, and the embeddings send $c_1$ to all
$\pm c_j$. Then $C_d=U_{d-2}(c_1)$ has the $2m$ conjugates $\pm g(\theta_j)$, which are distinct by Lemma~\ref{lem:gvalues}. Hence the degree is
$2m=d-1$.
\end{proof}
As a consistency check, exact minimal-polynomial computations for $3\le d\le12$ give precisely the degrees in Theorem~\ref{thm:degree};
for example, $C_4=9/7$ and $C_6=(225+140\sqrt{70})/961$ \cite[Cor.~5.1, Prop.~5.2]{I}.
The distinctness statement in Lemma~\ref{lem:gvalues} was also checked numerically for $4\le d\le400$.

\begin{corollary}\label{cor:degree}
$[\Q(C_d):\Q]$ equals $d-1$ (odd $d$) or $d/2-1$ (even $d$) for all $3\le d\le400$, for all even $d$ with $\gcd(v_2(d),d/2-1)=1$ or $d-1$ prime, and
for all $d=2p+1$ with $p$ prime and $\operatorname{ord}_p(2)=p-1$.
\end{corollary}

\section{Computations}\label{sec:comp}
All computations use Python 3.12.3, SymPy 1.14.0 and mpmath 1.3.0; the scripts and outputs are in the supplementary archive.

\subsection*{Irreducibility for $3\le d\le400$}
If $A_d=FG$ with $\deg F=k$, then for every prime $p$ not dividing the leading coefficient, $k$ is a sum of a sub-multiset of the degrees of the
irreducible factors of $A_d\bmod p$. We intersect these sets over primes until no $k\in[1,m/2]$ survives. The $398$ values split as follows:
$185$ by Theorem~\ref{thm:even}, $6$ by Proposition~\ref{prop:eis}, $22$ by Theorem~\ref{thm:2p1}, $2$ of degree $\le1$, and $183$ by such
certificates ($176$ odd $d$ and the exceptional $d=56,248,296,320,344,352,392$). All primes used are $\le53$. Exact factorization over $\Z$ for
$d\le120$ agrees.

\subsection*{Galois groups}
For $17\le d\le82$, $\Gal(A_d)\cong S_m$, certified by Frobenius elements at unramified primes: an $m$-cycle, a $(m-1,1)$-element, an element with
a cycle of prime length in $(m/2,m-3]$, and an odd element (Jordan's theorem). For $d=2p+1$ and all primes $3\le p\le199$, $\Gal\cong S_p$:
irreducibility holds by the above, Theorem~\ref{thm:Gstar} applies for $p\ge5$, and for $p\equiv7\pmod8$ an odd Frobenius is found at some
$q\in\{3,5,11\}$.

\subsection*{The family $d=2p+2$}
For every prime $5\le p\le3000$ a table records both local sides, with $\ell,a,k$, the exact residue $N_k\bmod\ell$, $E_\ell$, $M_\ell$, and the
dual $q,b,E_q,M_q$. It also records the certificate used and its exact margin. We find $424$ of $428$ primes certified by local blocks: $381$
by the dual block with Jordan--Marggraff, $31$ by the $\ell$-block with Jordan--Marggraff, and $12$ by a dual block with $b\ge2$ and (d). All
$424$ are free of CFSG. The minimal margin is $3/14$ for Jordan--Marggraff ($1/2-s/p$) and $1/4$ for (d) ($p/4-s$). The value $1/4$ occurs exactly when $p+1=2q^b$ (here $p=53,337$), since then $s=(p-1)/4$. The Liebeck--Saxl margin $1/3-s/p$
(Remark~\ref{rem:LS}) is at least $85/1011$. The four exceptions $p\in\{5,11,23,191\}$ are the primes
with $p=3\cdot2^s-1$ and $2p+1=3\cdot2^{s+1}-1$ both prime. For them $S_p$ is certified directly by Frobenius elements.

Parity was settled for every row: Theorem~\ref{thm:T2} in $216$ cases, Proposition~\ref{prop:PD} in $161$, Proposition~\ref{prop:PL} in $34$, and a
Stickelberger witness in $13$. A Stickelberger witness is a prime $q$ with $(\disc A_d/q)=-1$, which by Stickelberger--Swan \cite{Swan} gives an odd
Frobenius. Here $\disc A_d\bmod q$ is computed as a resultant over $\F_q$, cross-checked against exact discriminants in $366$ cases. Hence
$\Gal(A_{2p+2})\cong S_p$ for all primes $3\le p\le3000$ ($p=3$: a cubic with non-square discriminant).

For $p\le10^6$, Theorem~\ref{thm:Geven}(1) covers $78482$ of $78496$ primes. The other $14$ are exactly those with $p+1\in\{2^s,3\cdot2^s\}$:
nine are covered by Theorem~\ref{thm:Geven}(2), four are the exceptions above, and $p=524287$ was not checked. The local parity results alone give
$S_p$ for $98.2\%$ of these $p$.

\section{Open problems}\label{sec:open}
\begin{enumerate}
\item Conjecture~\ref{conj:main}; in particular the \emph{macroscopic} range of factor degrees for odd $d$. Example~\ref{ex:47} shows that a global
argument is needed for odd prime $d$.
\item The exceptional even set $\mathcal E$: find a uniform argument excluding the degrees $Q\,s$ allowed by Proposition~\ref{prop:ore}
(for $g=3$ only $Q$ and $2Q$). Remark~\ref{rem:Estruct} shows that local data at the structured primes do so for most, but not all, $d\in\mathcal E$.
\item Prove $\Gal(A_d)\supseteq A_m$ for composite $m$. Primitivity is the difficulty; the local cycles of Section~\ref{sec:blocks} restrict
block sizes but not the large ones.
\item Closed-form square classes for $a\ge2$, for unclean $\ell$, and for $q^b$ with $b\ge2$.
\item Related families: $U_d(2,t)-c$ for integers $c\ne0$ (irreducible for $c\in\{\pm1,\pm2,3,4\}$, $d\le50$), and the cosine analogue
$(V_d(2,t)-2)/(t-1)$ with $V_d=\alpha^d+\beta^d$ (the equation $\cos d\theta=\cos^d\theta$; irreducible for $d\le50$). In all neighbouring
families we tested, reducibility comes only from three mechanisms: a degenerate point $\beta=0$, $\alpha/\beta$ a root of unity, or
divisibility $U_k\mid U_d$.
\end{enumerate}

\section*{Use of AI tools}
AI systems used in the research and manuscript workflow included OpenAI GPT-5.6 Sol and Anthropic Claude Opus 5.5. They were used for
mathematical exploration, proof development and critique, symbolic and numerical verification, literature search, and manuscript drafting
and editing. AI systems are not authors. The author reviewed the mathematical statements and proofs and takes full responsibility for the
content.

\section*{Supplementary material}
Scripts, certificate tables and their checksums are provided in the supplementary archive accompanying this manuscript.

\begin{thebibliography}{99}
\bibitem{A} A.~Martinevski, \emph{Two-point extremals for characteristic-function dependence with fixed marginals}, preprint (2026).
\bibitem{Bochert97} A.~Bochert, \emph{Ueber die Classe der transitiven Substitutionengruppen II}, Math. Ann. \textbf{49} (1897), 133--144.
\bibitem{BHV} Yu.~Bilu, G.~Hanrot, P.~M.~Voutier, \emph{Existence of primitive divisors of Lucas and Lehmer numbers}, J. reine angew. Math.
\textbf{539} (2001), 75--122.
\bibitem{BL} Yu.~Bilu, F.~Luca, \emph{Binary polynomial power sums vanishing at roots of unity}, Acta Arith. \textbf{198} (2021), 195--217.
\bibitem{DKS} K.~Dilcher, S.-H.~Kim, K.~B.~Stolarsky, \emph{Properties of two Chebyshev-like polynomial sequences}, arXiv:2607.16940 (2026).
\bibitem{Dumas} G.~Dumas, \emph{Sur quelques cas d'irr\'eductibilit\'e des polyn\^omes \`a coefficients rationnels}, J. Math. Pures Appl. (6)
\textbf{2} (1906), 191--258.
\bibitem{GMN} J.~Gu\`ardia, J.~Montes, E.~Nart, \emph{Newton polygons of higher order in algebraic number theory}, Trans. Amer. Math. Soc.
\textbf{364} (2012), 361--416.
\bibitem{Hooley} C.~Hooley, \emph{On Artin's conjecture}, J. reine angew. Math. \textbf{225} (1967), 209--220.
\bibitem{I} A.~Martinevski, \emph{A sharp universal profile for multivariate characteristic-function factorization}, preprint (2026).
\bibitem{Jones} G.~A.~Jones, \emph{Primitive permutation groups containing a cycle}, Bull. Aust. Math. Soc. \textbf{89} (2014), 159--165.
\bibitem{LS} M.~W.~Liebeck, J.~Saxl, \emph{Minimal degrees of primitive permutation groups, with an application to monodromy groups of covers of
Riemann surfaces}, Proc. London Math. Soc. (3) \textbf{63} (1991), 266--314.
\bibitem{LT} R.~Levingston, D.~E.~Taylor, \emph{The theorem of Marggraff on primitive permutation groups which contain a cycle}, Bull. Austral. Math.
Soc. \textbf{15} (1976), 125--128.
\bibitem{Muller} P.~M\"uller, \emph{Permutation groups of prime degree, a quick proof of Burnside's theorem}, Arch. Math. \textbf{85} (2005), 15--17.
\bibitem{Ore} \O.~Ore, \emph{Newtonsche Polygone in der Theorie der algebraischen K\"orper}, Math. Ann. \textbf{99} (1928), 84--117.
\bibitem{Schinzel} A.~Schinzel, \emph{Primitive divisors of the expression $A^n-B^n$ in algebraic number fields}, J. reine angew. Math.
\textbf{268/269} (1974), 27--33.
\bibitem{Schomburg} B.~Schomburg, \emph{Bochert's results on the minimal degree of multiply transitive permutation groups}, arXiv:2209.12049 (2022).
\bibitem{Serre} J.-P.~Serre, \emph{Local Fields}, Graduate Texts in Mathematics \textbf{67}, Springer, 1979.
\bibitem{Swan} R.~G.~Swan, \emph{Factorization of polynomials over finite fields}, Pacific J. Math. \textbf{12} (1962), 1099--1106.
\bibitem{Wielandt} H.~Wielandt, \emph{Finite Permutation Groups}, Academic Press, 1964.
\end{thebibliography}
\end{document}
